\documentclass{article}
\usepackage{enumitem}
\usepackage{amsmath}
\begin{document}

\title{Chapter 6: Anatomy of Flowering Plants}
\date{}
\maketitle

\begin{enumerate}[label=Q\arabic*.]

\item The tissue system in plants that provides mechanical support and protection is:
\\(A) Vascular tissue system
\\(B) Ground tissue system
\\(C) Epidermal tissue system
\\(D) Meristematic tissue system

\item Which of the following is a simple permanent tissue?
\\(A) Xylem
\\(B) Phloem
\\(C) Parenchyma
\\(D) Epidermis

\item Sclereids (stone cells) are a type of:
\\(A) Parenchyma
\\(B) Collenchyma
\\(C) Sclerenchyma
\\(D) Xylem

\item The Casparian strip is present in the cell walls of:
\\(A) Epidermis of roots
\\(B) Endodermis of roots
\\(C) Cortex of stems
\\(D) Mesophyll of leaves

\item In a dicot stem, the vascular bundles are arranged:
\\(A) In a ring with conjoint, collateral, and open bundles
\\(B) Scattered throughout the ground tissue
\\(C) In a single central bundle
\\(D) Only in the pith region

\item Assertion: Secondary growth occurs in monocot stems.\\
Reason: Monocot stems lack vascular cambium between xylem and phloem.
\\(A) Both true, Reason is correct explanation
\\(B) Both true, Reason is not correct explanation
\\(C) Assertion false, Reason true
\\(D) Assertion true, Reason false

\item The tissue responsible for the formation of cork in older dicot stems is:
\\(A) Vascular cambium
\\(B) Phellogen (cork cambium)
\\(C) Interfascicular cambium
\\(D) Fascicular cambium

\item Which of the following correctly describes bulliform cells?
\\(A) Large, vacuolated epidermal cells in grasses that cause leaf rolling under drought
\\(B) Thick-walled cells providing mechanical support in stems
\\(C) Guard cell type found in C4 plants
\\(D) Cells forming the bundle sheath in C4 leaves

\item Annual rings (growth rings) in woody dicots are formed due to:
\\(A) Alternate formation of spring wood and autumn wood by vascular cambium
\\(B) Deposition of secondary phloem in concentric rings
\\(C) Seasonal changes in parenchyma cell size
\\(D) Activity of cork cambium during dry and wet seasons

\item The vascular bundle in monocot stems differs from dicot stems in being:
\\(A) Conjoint, collateral, and open (with cambium)
\\(B) Conjoint, collateral, and closed (without cambium), scattered
\\(C) Radial and exarch, arranged in a ring
\\(D) Concentric with phloem surrounding xylem

\item Companion cells are associated with which tissue?
\\(A) Xylem vessels
\\(B) Sieve tube elements of phloem
\\(C) Tracheids
\\(D) Sclerenchyma fibres

\item Which type of meristem is responsible for increase in girth of a plant?
\\(A) Apical meristem
\\(B) Intercalary meristem
\\(C) Lateral meristem
\\(D) Wound meristem

\item The lenticel in woody stems functions as:
\\(A) Site of photosynthesis in bark
\\(B) A pore allowing gaseous exchange through the periderm
\\(C) A structural support in the secondary xylem
\\(D) A water-absorbing structure

\item Match the following tissues with their functions:\\
1. Collenchyma $\rightarrow$ (a) Flexible mechanical support in growing organs\\
2. Sclerenchyma $\rightarrow$ (b) Rigid mechanical support in mature organs\\
3. Parenchyma $\rightarrow$ (c) Storage and metabolic functions\\
4. Aerenchyma $\rightarrow$ (d) Buoyancy in aquatic plants
\\(A) 1-a, 2-b, 3-c, 4-d
\\(B) 1-b, 2-a, 3-d, 4-c
\\(C) 1-c, 2-d, 3-a, 4-b
\\(D) 1-d, 2-c, 3-b, 4-a

\item The heartwood (duramen) in old trees differs from sapwood (alburnum) because:
\\(A) Heartwood actively conducts water; sapwood does not
\\(B) Heartwood is non-functional, dark, dense with tannins; sapwood is functional and lighter
\\(C) Heartwood is formed by phloem; sapwood is formed by xylem
\\(D) Heartwood lacks tracheids; sapwood has both tracheids and vessels

\end{enumerate}
\end{document}
